\documentclass[11pt, a4paper]{article}

\usepackage{fontspec}
\setmainfont{Latin Modern Roman}
\setmonofont{DejaVu Sans Mono}[Scale=0.85]
\usepackage[margin=2.5cm]{geometry}
\usepackage{amsmath, amssymb, amsthm}
\usepackage{hyperref}
\usepackage[dvipsnames]{xcolor}
\usepackage{listings}
\usepackage{enumitem}
\usepackage{fancyvrb}
\usepackage{caption}
\captionsetup[figure]{name=Listing}
\newcommand{\lean}[1]{\ifmmode\text{\ttfamily\detokenize{#1}}\else{\ttfamily\detokenize{#1}}\fi}
\newcommand{\brk}{\discretionary{}{}{}}
\emergencystretch=5em

\definecolor{indigo}{rgb}{0.29, 0.0, 0.51}
\definecolor{teal}{rgb}{0.0, 0.5, 0.5}

\hypersetup{
    colorlinks=true,
    linkcolor=teal,
    citecolor=indigo,
    urlcolor=blue
}

% --- LEAN 4 SYNTAX HIGHLIGHTING ---
\definecolor{keywordcolor}{rgb}{0.13, 0.13, 0.6}
\definecolor{stringcolor}{rgb}{0.6, 0.13, 0.13}
\definecolor{commentcolor}{rgb}{0.25, 0.5, 0.35}
\definecolor{leanbackground}{rgb}{0.97, 0.98, 0.99}

\lstdefinelanguage{lean4}{
  keywords={def, theorem, lemma, structure, class, where, by, sorry, intros, intro, rfl, contradiction, cases, induction, open, namespace, end, import, have, show, match, with, decide, constructor, let, in, noncomputable, exact, apply, rw, simp, calc, ring, positivity, linarith, unfold},
  keywordstyle=\color{keywordcolor}\bfseries,
  ndkeywords={Nat, Real, Bool, Prop, Type, String, Int, Float, List, Complex, UpperHalfPlane, GLPos2Real, SL2Z, ModularForm},
  ndkeywordstyle=\color{purple}\bfseries,
  comment=[l]{--},
  morecomment=[s]{/-}{-/},
  commentstyle=\color{commentcolor}\itshape,
  stringstyle=\color{stringcolor},
  sensitive=true,
  extendedchars=true,
  literate={ℝ}{{$\mathbb{R}$}}1 {ℂ}{{$\mathbb{C}$}}1 {ℤ}{{$\mathbb{Z}$}}1 {↑}{{$\uparrow$}}1 {≠}{{$\neq$}}1 {>}{{$>$}}1 {<}{{$<$}}1 {≥}{{$\geq$}}1 {≤}{{$\leq$}}1 {₀}{{$_0$}}1 {∀}{{$\forall$}}1 {⟨}{{$\langle$}}1 {⟩}{{$\rangle$}}1 {→}{{$\to$}}1
}

\lstset{
  language=lean4,
  backgroundcolor=\color{leanbackground},
  basicstyle=\ttfamily\small,
  breaklines=true,
  frame=single,
  rulecolor=\color{gray!30},
  captionpos=b,
  keepspaces=true,
  showstringspaces=false
}

\newtheorem{theorem}{Theorem}[section]
\newtheorem{lemma}[theorem]{Lemma}
\newtheorem{definition}[theorem]{Definition}
\theoremstyle{remark}
\newtheorem{remark}[theorem]{Remark}

\title{\textbf{A Mathlib-Idiomatic Formalization of the Fricke \\ Involution and its Operator on Modular Forms}\\[0.8em]
  \large\textcolor{BrickRed}{PREPRINT --- not peer reviewed}}
\author{
  Xavier Callens\thanks{Corresponding author: \texttt{callensxavier@gmail.com}} \\
  \small SocrateAI Research Initiative
}
\date{\today\\[0.6em]
  \small Archived at Zenodo: \href{https://doi.org/10.5281/zenodo.22542571}{\texttt{10.5281/zenodo.22542571}} (all versions)}

\begin{document}
\maketitle

\begin{center}
\fbox{\begin{minipage}{0.92\textwidth}
\small \textbf{Status.} This is a \emph{preprint}. It has \textbf{not been peer reviewed}. The
Lean development it describes compiles and is axiom-audited --- those claims are machine-checked
and independently reproducible --- but the \emph{exposition} has had no external referee. This
version incorporates revisions from two internal adversarial passes: one against an earlier claim
of mathematical priority (found false and retracted, \S\ref{sec:related}) and one against an
attempted T-duality bridge theorem (found unproved and reported as a negative result,
\S\ref{sec:tduality}); neither constitutes venue peer review. Read the scope before citing the
result.
\end{minipage}}
\end{center}

\begin{abstract}
We formalize, in Lean 4 on top of Mathlib, the Fricke involution
$W_N = \left(\begin{smallmatrix} 0 & -1 \\ N & 0 \end{smallmatrix}\right)$ on $\Gamma_0(N)$
and the operator it induces on modular forms, as a small reusable API rather than as
proof-internal infrastructure. We claim no priority: Fricke and Atkin--Lehner material already
exists in Lean~4 in the Fermat's Last Theorem artifact~\cite{FLT2026}
(\S\ref{sec:related}), which covers strictly more. What we offer is a packaging aimed at
upstreaming --- Mathlib itself contains no Fricke or Atkin--Lehner operator. Mathlib already provides the upper
half-plane, the M\"obius action of $GL_2(\mathbb{R})$, the weight-$k$ slash action carrying
the determinant factor, the congruence subgroups $\Gamma_0(N)$, and the bundled types
\lean{SlashInvariantForm} and \lean{ModularForm}; it contains no Fricke or Atkin--Lehner
operator. We supply that missing piece in three layers: (i) the matrix level --- $W_N \in
GL_2^+(\mathbb{R})$ with $\det W_N = N$, the involution identity $W_N^2 = -N \cdot I$, the
explicit conjugation $W_N \gamma W_N^{-1} = \left(\begin{smallmatrix} d & -c/N \\ -Nb & a
\end{smallmatrix}\right)$, and the theorem that $W_N$ normalizes $\Gamma_0(N)$; (ii) the
slash level --- $f \mapsto f\vert_k W_N$ preserves $\Gamma_0(N)$-slash-invariance; and
(iii) the bundled level --- the operator is well defined on \lean{ModularForm}, its
holomorphy and boundedness at the cusps transported along $W_N$.
The development is 27 declarations in 291 lines; all 21 theorems compile with no \lean{sorry}
and depend on exactly Lean's three standard axioms \lean{propext}, \lean{Classical.choice} and
\lean{Quot.sound}, pinned by 16 build-failing \texttt{\#guard\_msgs} checks so that a drifted
footprint fails compilation. We make no claim of new mathematics: the results are classical,
due to Fricke and to Atkin--Lehner~\cite{AtkinLehner1970}. The contribution is the
mechanization, and its intended use is as infrastructure for formalizing eta-quotient
identities, where Fricke eigenvalues are the organizing invariant.
\end{abstract}

\noindent\textbf{MSC 2020:} 11F11, 11F06, 68V20.\quad
\textbf{Keywords:} Fricke involution, Atkin--Lehner theory, congruence subgroups,
Lean 4, Mathlib, formal verification.

\section{Introduction}

The Fricke involution is elementary to state and ubiquitous in the theory of modular forms
on $\Gamma_0(N)$. For $N \ge 1$ set
\[
  W_N \;=\; \begin{pmatrix} 0 & -1 \\ N & 0 \end{pmatrix}, \qquad \det W_N = N > 0 .
\]
Acting on the upper half-plane by $\tau \mapsto -1/(N\tau)$, it exchanges the cusps $0$ and
$\infty$ of $\Gamma_0(N)$, normalizes $\Gamma_0(N)$, and induces an involution on
$M_k(\Gamma_0(N))$ whose $\pm 1$ eigenspaces organize much of the arithmetic of the space
--- most visibly for eta-quotients, where multiplicativity is read off from Fricke and
Atkin--Lehner eigenvalues~\cite{Martin1996, Ono2004}.

This note is written in the mode that recent practice makes standard: a formal artifact
shipped alongside the human exposition. Large-scale autoformalization has crossed a
threshold --- a complete machine-checked proof of Fermat's Last Theorem was produced in
eleven days by AI agents coordinating through a shared statement graph~\cite{FLT2026,
Prove2Me2026} --- and the surrounding methodological questions are now mainstream
mathematics~\cite{Tao2026AI, Buzzard2021}. We work in Lean~4~\cite{deMoura2021} against
Mathlib~\cite{Mathlib2020}, and we adopt two specific practices from the FLT artifact:
the axiom audit as a \emph{build target}, and an explicit statement-dependency graph
(\S\ref{sec:verification}).

\subsection*{What Mathlib already provides}

We are explicit about prior art, because it determines what is and is not new here.
As of commit \texttt{905b9581} (2026-07-28), Mathlib~\cite{Mathlib2020} contains, in forms
strictly more general than we would need to restate:

\begin{itemize}
  \item \lean{UpperHalfPlane} and the M\"obius action of \lean{GL (Fin 2) ℝ},
        together with \lean{denom_ne_zero_of_im} and \lean{normSq_denom_pos};
  \item \lean{im_smul_eq_div_normSq}, which is exactly the classical identity
        $\operatorname{Im}(M \cdot \tau) = |\det M| \operatorname{Im}(\tau) / |c\tau + d|^2$,
        stated for arbitrary $M \in GL_2(\mathbb{R})$ via $|\det M|$ rather than only for
        $\det M > 0$;
  \item \lean{Matrix.GLPos}, the subgroup of positive determinant;
  \item \lean{CongruenceSubgroup.Gamma0 N};
  \item the weight-$k$ slash action
        \lean{SlashAction ℤ (GL (Fin 2) ℝ) (ℍ → ℂ)}, whose definition
        \emph{already carries} the determinant factor,
        $f\vert_k g = \sigma(g)\, f(g\cdot\tau)\, |\det g|^{k-1} \operatorname{denom}(g,\tau)^{-k}$;
  \item the bundled \lean{ModularForm Γ k} for
        \lean{Γ : Subgroup (GL (Fin 2) ℝ)}, with holomorphy and boundedness at cusps.
\end{itemize}

None of that is reproved here; we import it. A reader looking for a formalization of the
$GL_2^+(\mathbb{R})$-action on $\mathbb{H}$ should use Mathlib's, not ours.

\subsection*{The gap, and what we contribute}

A repository-wide search of Mathlib at that commit for \texttt{Fricke} or
\texttt{AtkinLehner} returns \emph{zero} occurrences: the machinery to host these operators
exists in Mathlib, but the operators do not. That is the gap we address --- with the important
qualification, developed in \S\ref{sec:related}, that Fricke material \emph{does} already exist
in Lean~4 outside Mathlib. Our contributions are therefore stated as an API, not as firsts:

\begin{enumerate}
  \item[(C1)] $W_N$ as an element of \lean{GL (Fin 2) ℝ} with $\det W_N = N$, and
        $W_N \in GL_2^+(\mathbb{R})$ (Theorem~\ref{thm:glpos});
  \item[(C2)] $W_N^2 = -N \cdot I$ (Theorem~\ref{thm:involution}), which since scalar
        matrices act trivially on $\mathbb{H}$ is the involution property;
  \item[(C3)] the explicit conjugation formula (Theorem~\ref{thm:conj});
  \item[(C4)] $W_N$ conjugates $\Gamma_0(N)$ into itself inside $GL_2(\mathbb{R})$
        (Theorem~\ref{thm:normalizes}); together with $W_N^2 = -N\cdot I$ central this gives
        the two-sided normalizer condition, though we formalize the containment;
  \item[(C5)] slashing by $W_N$ preserves $\Gamma_0(N)$-slash-invariance, giving the operator
        on \lean{SlashInvariantForm} (Theorem~\ref{thm:slashinv});
  \item[(C6)] holomorphy and boundedness at the cusps transport along $W_N$, giving the
        operator on the bundled \lean{ModularForm} (Theorem~\ref{thm:modular});
  \item[(C7)] the composite $f \mapsto (f\vert_k W_N)\vert_k W_N$ is multiplication by the
        explicit scalar $(-1)^k N^{k-2}$ (Theorem~\ref{thm:composite}), which identifies the
        normalization making $\vert_k W_N$ an involution.
\end{enumerate}

(C4) is the hinge: it is what makes $f \mapsto f\vert_k W_N$ well defined at all. (C5) and
(C6) carry it up to Mathlib's bundled types, so that the Fricke operator can be composed with
the rest of the library rather than living beside it.

\section{Mathematical background}

Let $\gamma = \left(\begin{smallmatrix} a & b \\ c & d\end{smallmatrix}\right) \in SL_2(\mathbb{Z})$
with $\gamma \in \Gamma_0(N)$, i.e.\ $N \mid c$; write $c = Nc'$. A direct computation gives
\begin{equation}\label{eq:conj}
  W_N \gamma W_N^{-1}
  \;=\; \begin{pmatrix} d & -c' \\ -Nb & a \end{pmatrix} ,
\end{equation}
whose entries are integers precisely because $N \mid c$, whose determinant is
$da - (-c')(-Nb) = ad - bc = 1$, and whose lower-left entry $-Nb$ is divisible by $N$.
Hence the conjugate again lies in $\Gamma_0(N)$, and $W_N$ normalizes $\Gamma_0(N)$.

Rather than manipulate $W_N^{-1}$ in the formalization, we prove the equivalent
\emph{intertwining} identity
\begin{equation}\label{eq:intertwine}
  W_N \cdot \gamma \;=\; \delta \cdot W_N, \qquad
  \delta := \begin{pmatrix} d & -c' \\ -Nb & a \end{pmatrix},
\end{equation}
which avoids inverses entirely and from which \eqref{eq:conj} follows by one right
multiplication. Both sides of \eqref{eq:intertwine} equal
$\left(\begin{smallmatrix} -c & -d \\ Na & Nb \end{smallmatrix}\right)$, the two agreeing
exactly on the relation $c = Nc'$.

\section{The Lean development}

The development is four modules. We reproduce the load-bearing declarations verbatim from
the compiled source.

\begin{figure}[h]
\begin{Verbatim}[frame=single,fontsize=\small,samepage=true]
def frickeMatrix (N : ℕ) : Matrix (Fin 2) (Fin 2) ℝ := !![0, -1; (N : ℝ), 0]

@[simp] theorem frickeMatrix_det (N : ℕ) : (frickeMatrix N).det = (N : ℝ) := by
  simp [frickeMatrix, Matrix.det_fin_two_of]

noncomputable def frickeW {N : ℕ} (hN : 0 < N) : GL (Fin 2) ℝ :=
  Matrix.GeneralLinearGroup.mkOfDetNeZero (frickeMatrix N) (frickeMatrix_det_ne_zero hN)
\end{Verbatim}
\caption{The Fricke matrix and its determinant.}
\end{figure}


\begin{theorem}[$W_N$ has positive determinant]\label{thm:glpos}
\lean{frickeW_mem_GLPos : frickeW hN ∈ GLPos (Fin 2) ℝ}.
\end{theorem}

\begin{theorem}[Involution]\label{thm:involution}
$W_N^2 = -N \cdot I$:
\lean{frickeW_sq_coe : ((frickeW hN * frickeW hN : GL (Fin 2) ℝ) : Matrix _ _ ℝ)}
\lean{= (-(N : ℝ)) • 1}.
\end{theorem}

Since scalar matrices act trivially on $\mathbb{H}$, Theorem~\ref{thm:involution} says
precisely that $W_N$ induces an involution of the upper half-plane.

The arithmetic input is the divisibility extracted from membership in $\Gamma_0(N)$, which
Mathlib states modulo $N$:

\begin{figure}[h]
\begin{Verbatim}[frame=single,fontsize=\small,samepage=true]
theorem gamma0_dvd_lower_left {N : ℕ} {γ : SL(2, ℤ)} (hγ : γ ∈ Gamma0 N) :
    (N : ℤ) ∣ γ.1 1 0 := by
  rw [Gamma0_mem] at hγ
  exact (ZMod.intCast_zmod_eq_zero_iff_dvd _ _).mp hγ
\end{Verbatim}
\caption{Divisibility of the lower-left entry.}
\end{figure}


The conjugate $\delta$ of \eqref{eq:conj} is then definable as a genuine $SL_2(\mathbb{Z})$
element, its determinant condition discharged by \lean{linear_combination} against
$\det \gamma = 1$:

\begin{figure}[h]
\begin{Verbatim}[frame=single,fontsize=\footnotesize,samepage=true]
def frickeConj {N : ℕ} (γ : SL(2, ℤ)) (c' : ℤ) (hc : γ.1 1 0 = (N : ℤ) * c') : SL(2, ℤ) :=
  ⟨!![γ.1 1 1, -c'; -(N : ℤ) * γ.1 0 1, γ.1 0 0], by
    have hdet := γ.2
    rw [Matrix.det_fin_two] at hdet; rw [hc] at hdet
    rw [Matrix.det_fin_two_of]; linear_combination hdet⟩
\end{Verbatim}
\caption{The Fricke conjugate, with its determinant proof.}
\end{figure}


\begin{theorem}[Conjugation]\label{thm:conj}
\lean{frickeW_conj_eq} : for $\gamma \in SL_2(\mathbb{Z})$ with $c = Nc'$,
\[
  W_N \cdot \mathrm{mapGL}_{\mathbb{R}}(\gamma) \cdot W_N^{-1}
  \;=\; \mathrm{mapGL}_{\mathbb{R}}(\mathrm{frickeConj}\,\gamma\,c') .
\]
The proof rewrites by the intertwining identity \eqref{eq:intertwine} and cancels
$W_N W_N^{-1}$.
\end{theorem}

\begin{theorem}[Main result: normalization]\label{thm:normalizes}
\lean{frickeW_normalizes_Gamma0} :
\[
  \forall g \in \mathrm{Gamma0}(N).\mathrm{map}(\mathrm{mapGL}_{\mathbb{R}}), \quad
  W_N \, g \, W_N^{-1} \in \mathrm{Gamma0}(N).\mathrm{map}(\mathrm{mapGL}_{\mathbb{R}}) .
\]
\end{theorem}

\begin{proof}[Proof (as formalized)]
Destructure $g = \mathrm{mapGL}_{\mathbb{R}}(\gamma)$ with $\gamma \in \Gamma_0(N)$;
obtain $c'$ with $c = Nc'$ from \lean{gamma0_dvd_lower_left}; rewrite by
Theorem~\ref{thm:conj}; conclude by \lean{Subgroup.mem_map_of_mem} applied to
\lean{frickeConj_mem_Gamma0}, whose content is that the lower-left entry $-Nb$
vanishes modulo $N$.
\end{proof}

\section{The operator on modular forms}

Normalization is a statement about groups; being an \emph{operator on modular forms} requires
carrying the two analytic side conditions of \lean{ModularForm} along $W_N$. We do this in two
steps, working with the image $\Gamma_0(N)_{\mathbb{R}} := \Gamma_0(N).\mathrm{map}(\mathrm{mapGL}_{\mathbb{R}})$
inside $GL_2(\mathbb{R})$, which is the subgroup Mathlib's bundled types are stated over.

\begin{theorem}[Slash-invariance]\label{thm:slashinv}
\lean{slash_frickeW_invariant} : for $f$ a \lean{SlashInvariantForm} of weight $k$ for
$\Gamma_0(N)_{\mathbb{R}}$ and any $g$ in that subgroup,
$(f\vert_k W_N)\vert_k g = f\vert_k W_N$.
\end{theorem}

\begin{proof}[Proof (as formalized)]
Write $g = \mathrm{mapGL}_{\mathbb{R}}(\gamma)$ with $\gamma \in \Gamma_0(N)$ and extract
$c = Nc'$. Then, using the slash cocycle in both directions and the intertwining
identity~\eqref{eq:intertwine},
\[
  (f\vert_k W_N)\vert_k \gamma \;=\; f\vert_k(W_N\gamma) \;=\; f\vert_k(\delta W_N)
  \;=\; (f\vert_k \delta)\vert_k W_N \;=\; f\vert_k W_N ,
\]
the last step by slash-invariance of $f$ under $\delta \in \Gamma_0(N)$
(Theorem~\ref{thm:conj}). This packages as \lean{frickeSlashOperator}.
\end{proof}

\begin{theorem}[Operator on \lean{ModularForm}]\label{thm:modular}
\lean{frickeModularOperator} : $f \mapsto f\vert_k W_N$ is a map
$\lean{ModularForm}\ \Gamma_0(N)_{\mathbb{R}}\ k \to \lean{ModularForm}\ \Gamma_0(N)_{\mathbb{R}}\ k$.
\end{theorem}

\begin{proof}[Proof (as formalized)]
Holomorphy is immediate from Mathlib's \lean{MDifferentiable.slash}: slashing by any fixed
element of $GL_2(\mathbb{R})$ preserves \lean{MDiff}. For boundedness, Mathlib's
\lean{OnePoint.IsBoundedAt.smul_iff} reduces \emph{$f\vert_k W_N$ is bounded at $c$} to
\emph{$f$ is bounded at $W_N \cdot c$}, so it suffices that $W_N$ carries cusps of
$\Gamma_0(N)$ to cusps of $\Gamma_0(N)$. That is \lean{isCusp_frickeW_smul}, obtained from
Mathlib's \lean{IsCusp.smul} together with the subgroup form of Theorem~\ref{thm:normalizes}
(\lean{frickeW_conjAct_le}) and the monotonicity of \lean{IsCusp} in its subgroup argument.
\end{proof}

Concretely $W_N$ exchanges the cusps $0$ and $\infty$ of $\Gamma_0(N)$; what is proved here is
the weaker statement, sufficient for the transport, that the cusp set is preserved.

\begin{theorem}[The composite is a scalar]\label{thm:composite}
\lean{frickeW_sq_slash} : for every $f : \mathbb{H} \to \mathbb{C}$ and every weight $k$,
\[
  (f\vert_k W_N)\vert_k W_N \;=\; (N^2)^{k-1}(-N)^{-k}\, f \;=\; (-1)^k N^{k-2}\, f .
\]
\end{theorem}

\begin{proof}[Proof (as formalized)]
By the cocycle law the left side is $f\vert_k(W_N^2)$, and $W_N^2 = -N\cdot I$ is scalar
(Theorem~\ref{thm:involution}). For that matrix the three ingredients of Mathlib's slash are
computed directly: $\det(W_N^2) = N^2$, $\operatorname{denom}(W_N^2,\tau) = -N$, and
$W_N^2$ acts trivially on $\mathbb{H}$. The conjugation factor $\sigma$ needs no determinant
case split: $\sigma(gg') = \sigma g \circ \sigma g'$ and $\sigma g \circ \sigma g = \mathrm{id}$
give $\sigma(W_N^2) = \mathrm{id}$ at once.
\end{proof}

For even $k$ it follows that $f \mapsto N^{1-k/2}(f\vert_k W_N)$ squares to the identity --- the
Fricke involution proper. For odd $k$ the question is vacuous, since $-I \in \Gamma_0(N)$ forces
$f \equiv 0$ (Mathlib's \lean{ModularForm.eq_zero_of_neg_one_mem}).

\section{Verification status}\label{sec:verification}

Every declaration in the module compiles. There are no \lean{sorry} placeholders and
no project-local axioms. The axiom footprint of each theorem, obtained by
\texttt{\#print axioms}, is exactly
\[
  [\,\texttt{propext},\ \texttt{Classical.choice},\ \texttt{Quot.sound}\,] ,
\]
the three standard axioms of Lean's logic. Following the Fermat's Last Theorem
artifact~\cite{FLT2026}, the footprint is not merely reported but \emph{enforced}: a
\lean{FinalCheck} module places a \texttt{\#guard\_msgs in \#print axioms} guard on each of the
sixteen headline theorems inside the default build target, so any widening of the axiom set is
a compile error, and a negative control (a guard asserting a deliberately wrong footprint)
is kept outside the build and verified to fail. The development is tracked as a
statement-dependency graph (\lean{dag/theorems.jsonl}, with a human-readable map in
\lean{dag/PROOF-PATH.md}); the machine checks that every node marked proved names an
existing declaration and depends only on proved nodes. We claim nothing beyond what the
kernel checks:
in the terminology we adopt throughout the SocrateAI programme, (C1)--(C7) are
\emph{formal} results; the surrounding exposition is not.

\paragraph{Artifact.}
Archived at Zenodo, \href{https://doi.org/10.5281/zenodo.22542571}{\texttt{doi:10.5281/zenodo.22542571}},
which is the citable version of record for this note and its sources.
Repository \url{https://github.com/xaviercallens/SocrateAI-Lean-Lib} (tag
\texttt{fricke-v1}); module
\texttt{Lean/\brk SocrateAI/\brk ModularForms/\brk \{FrickeInvolution,\brk\ FrickeSlash,\brk\ FrickeModular,\brk\ FrickeComposite\}.lean},
with the axiom guards in \texttt{Lean/\brk SocrateAI/\brk FinalCheck.lean};
a mirror with a dataset card is at
\url{https://huggingface.co/datasets/callensxavier/socrateai-fricke-involution}.
Toolchain \texttt{leanprover/lean4:v4.32.2};
Mathlib pinned at commit \texttt{905b9581} (2026-07-28);
\texttt{lake build SocrateAI} completes in 3053 jobs (including the axiom-guard module).
The build configuration is distributed and portable: \texttt{lakefile.lean} requires Mathlib from
git at the pinned revision above, \texttt{lean-toolchain} matches it, and no tracked file contains a
machine-specific path. \texttt{git clone}, \texttt{lake exe cache get}, \texttt{lake build SocrateAI}
reproduces the artifact; see \texttt{BUILDING.md}. Earlier versions of this note stated that the
configuration was not distributed. That was true, and understated: the published
\texttt{lakefile.lean} declared no Mathlib dependency at all and its \texttt{lean-toolchain} named a
different compiler version, so a clean clone could not build. Both are fixed.
The four modules total 27 declarations in 291 lines and import only
\texttt{Mathlib.\allowbreak NumberTheory.\allowbreak ModularForms.\allowbreak SlashActions},
\texttt{Mathlib.\allowbreak NumberTheory.\allowbreak ModularForms.\allowbreak CongruenceSubgroups},
\texttt{Mathlib.\allowbreak NumberTheory.\allowbreak ModularForms.\allowbreak SlashInvariantForms} and
\texttt{Mathlib.\allowbreak NumberTheory.\allowbreak ModularForms.\allowbreak Basic}.
A statement-dependency graph with a natural-language description per node, and a validator
enforcing that every node marked proved names an existing declaration and depends only on
proved nodes, is in \texttt{dag/}.

\section{Related work}\label{sec:related}

\paragraph{Prior Lean 4 formalization of Fricke material.}
The Fermat's Last Theorem artifact~\cite{FLT2026} is a Lean~4 development built on Mathlib that
already contains substantial Fricke and Atkin--Lehner theory. A repository-scoped code search
returns on the order of $2.7\times10^3$ Lean files mentioning \texttt{Fricke} and
$3.6\times10^3$ mentioning \texttt{AtkinLehner}, including
\texttt{Thm\_ModularForm\_heckeU\_add\brk\_slash\brk\_fricke\brk\_eq\brk\_zero},
\texttt{Thm\_CuspForm\_hasNebentypus\brk\_inv\brk\_and\brk\_qCoeff\brk\_hecke\brk\_eigen\brk\_of\brk\_fricke} and
\texttt{Def\_ModularCurve\_AtkinLehner}. Their statements use the same matrix
$\left(\begin{smallmatrix} 0 & -1 \\ N & 0\end{smallmatrix}\right)$ and the same Mathlib slash
action that we use here, and they go considerably further --- to Hecke operators, eigenforms with
nebentypus, and Atkin--Lehner operators on modular curves.

Most directly, \texttt{Def\_CuspForm\_AtkinLehnerOperator.lean} defines
\texttt{atkinLehnerLin}, a map
$\texttt{ModularForm}\ (\Gamma_0(M))\ k \to_{\ell[\mathbb{C}]} \texttt{ModularForm}\ (\Gamma_0(M))\ k$,
discharging slash-invariance, holomorphy and boundedness at the cusps as our
\texttt{frickeModularOperator} does --- our (C5) and (C6), for the whole Atkin--Lehner family and
as a $\mathbb{C}$-linear map. Their \texttt{Def\_ModularForm\_AtkinLehnerDatum.lean} builds its
group element with the same \texttt{mkOfDetNeZero} idiom we use and proves the analogue of our
(C1), (C2) and (C7).

In particular their \texttt{CohCarrier.frickeMat} is, up to naming, our
\texttt{frickeConj}: the same conjugate matrix
$\left(\begin{smallmatrix} d & -c/N \\ -Nb & a\end{smallmatrix}\right)$, the same divisibility
hypothesis, and a determinant proof that likewise finishes with
\texttt{linear\_combination}. Our development was arrived at independently, which does not bear
on priority; we record the coincidence because a reader comparing the two should know they are
the same construction.

\paragraph{What is nonetheless offered here.}
The FLT material is proof-internal: it lives in a bespoke \texttt{CohCarrier} namespace, and
$W_N$ is generally passed to each theorem as a parameter constrained by a matrix equation. We
instead define $W_N$ once as a \texttt{GL (Fin 2) ℝ} element and prove the small surrounding
theory --- membership in $GL_2^+$, $W_N^2 = -N\cdot I$, the conjugation formula, and the induced
operators on \texttt{SlashInvariantForm} and \texttt{ModularForm}. The intended destination is
Mathlib, which has none of this. Stated precisely, what is left after disclosure is: the literal
matrix $W_N$ over $\mathbb{R}$ (their datum's lower-right entry is $q$, so it does not specialize
to $W_N$ on the nose), the explicit composite scalar of (C7), and a four-named-import development
using stock tactics rather than a bespoke proof harness. That is a methodology difference, and
whether it is worth anything is a question for Mathlib review, not for us to assert.

\section{Limitations}

Fricke and Atkin--Lehner operators are classical; the definitive treatment of the
Atkin--Lehner involutions $W_Q$ for $Q \| N$ is~\cite{AtkinLehner1970}, and textbook
accounts appear in~\cite{DiamondShurman2005, Miyake2006, Ono2004}. Our contribution is not
mathematical novelty but the transfer of this material into a machine-checked setting where
it can be depended upon.

We state the limitations plainly.

\begin{remark}[What is \emph{not} done here]
\leavevmode
\begin{enumerate}[label=(\alph*)]
  \item We formalize the Fricke involution $W_N$ only. The full Atkin--Lehner family $W_Q$
        for $Q \| N$, and the resulting simultaneous eigenspace decomposition, are not
        treated.
  \item The $\pm 1$ eigenspace decomposition of $M_k(\Gamma_0(N))$ is not itself packaged as a
        Lean definition. The mathematical content is Theorem~\ref{thm:composite}, which is
        formalized; what remains is bookkeeping --- carrying the even-$k$ hypothesis through the
        \lean{zpow} arithmetic on the constant and bundling the two eigenspaces as submodules.
  \item We make no priority claim. Our gap check was run against Mathlib, where it holds; a
        later search of the Lean~4 developments cited here found the prior art described in
        \S\ref{sec:related}. We have not searched the Lean Zulip or other proof assistants.
\end{enumerate}
\end{remark}

\section{Intended application}

The motivation is downstream. Eta-quotients
$f(\tau) = \prod_{\delta \mid N} \eta(\delta \tau)^{r_\delta}$ are modular on $\Gamma_0(N)$
under Ligozat's conditions, and their behaviour under Fricke and Atkin--Lehner operators is
what makes multiplicativity tractable~\cite{Martin1996}. A formalization of eta-quotient
identities therefore needs $W_N$ as a prerequisite, which is what this module supplies.
A first attempt to make that link precise is reported in \S\ref{sec:tduality} below, and it
failed instructively rather than succeeding. With the operator now defined on \lean{ModularForm}, the immediate next steps are the
normalization constant of limitation (b) and the full Atkin--Lehner family of limitation (a). On the
analytic side, the Rademacher-type coefficient bounds that accompany the eta-quotient
programme will first be checked at the program-verified tier with lightweight
inequality-proof tooling in the style of Tao's \emph{estimates} assistant~\cite{TaoEstimates2025},
and only then promoted to kernel-checked statements --- the same two-tier discipline
(exact checker below, kernel above) used throughout this programme.

\section{A negative result: an attempted T-duality bridge}\label{sec:tduality}

A follow-up effort, reported here for the record rather than left as an unqualified forward
pointer, attempted to formalize that $W_N$ is not merely a device for eta-quotients but the
Lean~4 incarnation of a physical duality: that T-duality's action on the complex-structure
modulus $\tau$ of a $T^2$ factor, restricted to $\Gamma_0(N)$, \emph{coincides} with $W_N$'s
already-proven action on $\mathbb{H}$. This section states plainly what was found, because a
withheld negative result is exactly the failure mode this programme has already corrected twice
before~\cite{FLT2026} --- and because what survives the attempt is genuine and reusable.

\subsection*{What is proved, sorry-free, and reusable}

\begin{theorem}[$W_N$'s Möbius action, in closed form]\label{thm:tdualm1}
\lean{frickeW_smul_coe}: for every $N>0$ and $\tau\in\mathbb{H}$,
\[
  \bigl(W_N \cdot \tau : \mathbb{C}\bigr) \;=\; \frac{-1}{N\tau}.
\]
\end{theorem}

This closes a real seam that nothing in the library previously closed: the $-1/(N\tau)$ normal
form is exactly what the $\mathbb{C}$-valued eta-quotient half of this programme already computes
with (\lean{eta_fricke_factor}, in the eta-quotient modules), while $W_N$ itself had only ever
been stated in $GL_2(\mathbb{R})$-action language. The proof goes through Mathlib's determinant
case split (Theorem~\ref{thm:glpos} supplies exactly the positive-determinant hypothesis Mathlib's
own action lemma needs, so no hidden sign or conjugation enters) and is certified non-vacuous by a
negative control: the identical proof script against four independently-computed \emph{wrong}
right-hand sides ($+1/(N\tau)$, $-1/\tau$, $-N/\tau$, $-1/(N+\tau)$) fails with four distinct
\lean{unsolved goals}, each displaying the same correct left-hand side.

\subsection*{What is not proved, and why}

The bridge statement itself,
\[
  \bigl(W_N \cdot \tau : \mathbb{C}\bigr) = \tfrac{-1}{N\tau}
  \ \wedge\ \forall\,\gamma\in\Gamma_0(N),\ \gamma\cdot\tau
    \;\text{is the standard fractional-linear action of}\ SL_2(\mathbb{Z}),
\]
remains \lean{sorry}, under a build-failing inverted axiom guard, and an adversarial review of
this programme's own workflow found two independent reasons it should stay that way.

\textbf{It is a relabelling, not a reduction.} A standalone equivalence check, compiled and
verified, shows the bridge statement is logically equivalent to the conjunction of three already-
proved lemmas and adds no content beyond them: the hypothesis $\gamma\in\Gamma_0(N)$ is never
consumed (the second conjunct holds for every $\gamma\in SL_2(\mathbb{Z})$), and neither
Theorem~\ref{thm:normalizes} (that $W_N$ normalizes $\Gamma_0(N)$) nor the involution property
$W_N^2=-N\cdot I$ occurs anywhere in the statement or its proof term. Nothing in the Lean term
names a duality group, a string background, or a spectrum; the physics is asserted only in prose
around the theorem, not inside it.

\textbf{The physics citation does not support it.} The natural reference for a $\tau$-modulus
T-duality action, Giveon--Porrati--Rabinovici~\cite{GPR1994}, anchors $\tau\mapsto-1/\tau$ only at
$N=1$. For $N>1$, $\det W_N = N$, so $W_N$ is not even an element of the $SL_2(\mathbb{Z})$ that
reference describes, and an exhaustive search of its text returns no occurrence of a Fricke
involution, a congruence subgroup, or a level-dependent duality element at any $N>1$. The correct
physical home for a level-$N$ Fricke-type map turns out to be a different object entirely:
Persson--Volpato~\cite{PerssonVolpato2015} identify $S\mapsto -1/(NS)$ as an \emph{S-duality} on
the heterotic axio-dilaton modulus in CHL orbifold models, explicitly \emph{outside} the standard
$SL_2(\mathbb{Z})$ symmetry of the parent theory --- a different modulus, and a different duality,
than the one this attempt targeted.

\subsection*{The honest conclusion}

Theorem~\ref{thm:tdualm1} is real: it is evidence that $W_N$ is the right mathematical object for
a physical correspondence, stated purely as a fact about $\mathbb{H}$ and $\mathbb{C}$ with no
physics content claimed. Which correspondence, and under what hypotheses, remains open, and the
Persson--Volpato reference is where the next attempt should start --- inside CHL orbifolds of
$K3\times T^2$, on the axio-dilaton, not the $T^2$ modulus. We record this as a finding rather than
quietly dropping the sentence that used to promise it.

\section{Conclusion}

We have formalized the Fricke involution $W_N$ on $\Gamma_0(N)$ in Lean 4 and the operator it
induces on modular forms, addressing a gap in \emph{Mathlib} --- though not, as
\S\ref{sec:related} records, in Lean 4 as a whole. The development is
small --- 27 declarations in 291 lines --- depends only on Lean's three standard axioms, and is
built on Mathlib's existing \lean{SlashAction}, \lean{SlashInvariantForm} and
\lean{ModularForm} infrastructure rather than duplicating it, so that it composes with the
rest of the library. The mathematics is classical; the machine-checked statement of it is not.

\section*{A note on method: AI assistance and human responsibility}

This note follows the disclosure norm argued for by Tao~\cite{Tao2026AI}, endorsing recommendations
of the Leiden Declaration on Artificial Intelligence and Mathematics~\cite{Leiden2026}: disclose
tool use, affirm that credit and responsibility belong to the human author, and attribute sources
where AI tools cannot reliably do so themselves.

\textbf{What AI did.} The Lean~4 formalization, the literature search, and the drafting of this
text were produced with the assistance of large language models (Anthropic's Claude family),
working under human direction across many sessions. The methodology is neuro-symbolic in the
specific sense Tao describes~\cite{Tao2026AI}: natural-language reasoning proposes definitions,
proof strategies, and prose, and the Lean kernel --- not the model, and not the author's trust in
the model --- is what admits or rejects every mathematical claim. Every theorem in this paper
compiles against Mathlib with no \lean{sorry} and a build-failing axiom guard on its footprint
(\S\ref{sec:verification}); nothing here is asserted true because a model said so.

\textbf{What the author did.} Xavier Callens set the research direction, chose what to formalize
and in what order, judged what the results mean and whether they are worth publishing, and is
responsible for every claim in this document, including its errors. That responsibility was
exercised concretely, not nominally, in this paper's own history: an earlier draft claimed priority
for this formalization, and that claim was found false, retracted, and replaced with the
attribution recorded in \S\ref{sec:related}; a follow-up attempt to extend $W_N$ to a physical
duality (\S\ref{sec:tduality}) was found, on adversarial review, to reduce to a relabelling with an
unsupported physics citation, and is reported here as a negative result rather than closed as a
theorem. Tao's own test is apt: a result that cannot be explained and defended by its author,
formally verified or not, is not ready to be called complete~\cite{Tao2026AI}. The mathematical
content of $W_N$ --- what it is, why it normalizes $\Gamma_0(N)$, and why the two withdrawn claims
above were in fact wrong --- is something the author can explain without further AI assistance;
the Lean proof terms that certify it were substantially AI-drafted and are trusted only because the
kernel, not the author, checks them.

\begin{thebibliography}{99}
\bibitem{AtkinLehner1970} A.~O.~L.~Atkin, J.~Lehner, \textit{Hecke operators on $\Gamma_0(m)$}, Math.\ Ann.\ \textbf{185} (1970), 134--160. \texttt{doi:10.1007/BF01359701}.
\bibitem{Mathlib2020} The mathlib Community, \textit{The Lean mathematical library}, Proceedings of the 9th ACM SIGPLAN International Conference on Certified Programs and Proofs (CPP), 2020.
\bibitem{Buzzard2021} K.~Buzzard, \textit{Formalising mathematics}, Math.\ Proc.\ Cambridge Philos.\ Soc., 2021.
\bibitem{DiamondShurman2005} F.~Diamond, J.~Shurman, \textit{A First Course in Modular Forms}, Graduate Texts in Mathematics \textbf{228}, Springer (2005).
\bibitem{Miyake2006} T.~Miyake, \textit{Modular Forms}, Springer Monographs in Mathematics (2006).
\bibitem{Ono2004} K.~Ono, \textit{The Web of Modularity: Arithmetic of the Coefficients of Modular Forms and $q$-series}, CBMS Regional Conference Series in Mathematics \textbf{102}, AMS (2004). \texttt{doi:10.1090/cbms/102}.
\bibitem{Martin1996} Y.~Martin, \textit{Multiplicative eta-quotients}, Trans.\ Amer.\ Math.\ Soc.\ \textbf{348} (1996), 4825--4856.
\bibitem{GPR1994} A.~Giveon, M.~Porrati, E.~Rabinovici, \textit{Target Space Duality in String Theory}, Phys.\ Rept.\ \textbf{244} (1994), 77--202. \texttt{arXiv:hep-th/9401139}.
\bibitem{PerssonVolpato2015} D.~Persson, R.~Volpato, \textit{Fricke S-duality in CHL models}, JHEP \textbf{12} (2015) 156. \texttt{arXiv:1504.07260}.
\bibitem{FLT2026} Anthropic, \textit{A formal proof of Fermat's Last Theorem in Lean 4}, software artifact, \url{https://github.com/anthropics/fermats-last-theorem}, 2026.
\bibitem{Prove2Me2026} S.~Chen, K.~Marwaha, X.~Lu, H.~Yuen, T.~Peng, \textit{Prove2Me: An open collaborative platform for scaling math formalization}, arXiv:2608.28433, 2026.
\bibitem{Tao2026AI} T.~Tao, \textit{Mathematics in the age of AI}, essay based on a lecture at the
International Congress of Mathematicians, July 2026, submitted to the Proceedings of the ICM 2026.
\texttt{arXiv:2608.16753}.
\bibitem{Leiden2026} \textit{The Leiden Declaration on Artificial Intelligence and Mathematics},
June 2, 2026, endorsed by the International Mathematical Union. \url{https://leidendeclaration.ai/},
\texttt{doi:10.5281/zenodo.20302944}.
\bibitem{TaoEstimates2025} T.~Tao, \textit{estimates: a lightweight proof assistant for inequalities and asymptotic estimates}, software, \url{https://github.com/teorth/estimates}, 2025.
\bibitem{deMoura2021} L.~de~Moura et al., \textit{The Lean 4 Theorem Prover and Programming Language}, CADE-28, LNCS \textbf{12699}, Springer (2021).
\end{thebibliography}

\end{document}
